% Usability = 
% - Effectiveness: Can users actually achieve their goals with the system?
% - Efficiency: How quickly and with how much effort can users achieve their goals? This includes factors like time taken, number of steps, and cognitive load.
% - Satisfaction: Do users find the system enjoyable and pleasant to use? Are they satisfied with the overall experience?

% Ergebnisse interpretieren
% in den Forschungskontext einordnen


Systematischer Aufbau:
- Forschungsfrage(n) klar definiert
- Ergebnisse anhand einer bestimmten Methodik oder eines Forschungsprozesses interpretiert werden können
Mögliche Gliederungspunkte sind:
Beantwortung der Forschungsfragen
Interpretation der Ergebnisse
Diskussion der Limitationen der Forschung
Vorschläge für zukünftige Forschung

Vorteil:
- Fokussierung auf den Forschungsprozess und die einzelnen Schritte
- klare Darstellung der Ergebnisse im Kontext der Fragestellung
%Nachteil: Kann zu einer weniger inhaltlichen Gliederung führen, wenn die einzelnen Schritte nicht klar voneinander abgegrenzt sind.
- anhand der Überprüfung dieser Hypothesen strukturieren
-> wobei jeder Abschnitt die Ergebnisse zu einer bestimmten Hypothese diskutiert 

Weitere Aspekte:

    Es ist wichtig, dass die Diskussion nicht nur die Ergebnisse zusammenfasst, sondern auch interpretiert und kritisch hinterfragt.
    Die Ergebnisse sollten mit der vorhandenen Literatur verglichen und im Kontext des Forschungsstandes diskutiert werden.
    Die Limitationen der eigenen Forschung sollten ebenfalls erwähnt und ihre Auswirkungen auf die Ergebnisse diskutiert werden.
    Schließlich sollten Vorschläge für zukünftige Forschung gemacht werden, die auf den Ergebnissen der eigenen Arbeit basieren. 

-> Zusammenfassung Hypothesen neu bestätigen / ablehnen

Diskussion
->	nur das „warum ist das so“

- Was bedeuten die Unterschiede?
- Warum könnten sie so ausgefallen sein?
- Wie passen sie zu deiner Hypothese oder zu früheren Studien?
- Welche Limationen gibt es?
- Was könnte man daaus für die Praxis ableiten?

Ziel: Leser sollen verstehen, was die Ergebnisse bedeuten

\noindent
Die Hypothesen die nicht bestätigt werden können\\
------------------------------------\\
gut für mich\\
--> können nicht bestätigte Hypothesen durch ungünstige Fälle doch bestätigt werden?\\
--> wird das Ergebnis noch signifikanter?\\
schlecht für mich\\
--> muss eine Gegendarstellung betrachtet werden?\\
--> wo führen Extrema zur Signifikanz, die sonst nicht da ist?\\